\section{Turn 3 and the precise route-selection result}
The canonical nonresidue-factor graph and its two source character laws are
inherited from Bryan \cite[\S\S1--7]{BryanGraph}. The primitive sextic local form
and consecutive-defect examples are also antecedents
\cite[\S\S3--8]{BryanSix}\cite[\S\S2--5]{BryanChain}. Turn 1 retained the actual
stabilizer and inactive saturation, and Turn 2 constructed closed failed components
whose effective indices vary \cite{Turn1,Turn2}.

Here the stronger proposed assertion, that a cycle cannot remain primitive sextic
at every vertex, is disproved by a closed triangle at the prime
\[
 \boxed{p=7510085481569082811681\equiv121\pmod{840}.}       \tag{1}
\]
The complete original arithmetic, including primality, is certified below.
Cubic reciprocity and the actual cubic nilpotent component are consistent with the
triangle. A different result survives: a separately marked square multiplier
forces a nonresidue factor outside every canonical all-$3\bmod4$ triangle.
More generally, small primitive closed components cannot remain closed after the
square multipliers $1,9,25$ are all admitted. This is escape from a specified
component, not a universal theorem that a later selector is occupied.

\section{Original source, targets, and marked return}
Let $p$ be a prime in one of the six classes $1,121,169,289,361,529\pmod{840}$.
For a prime $q<p$ with $(q/p)=-1$, the canonical multiplier is
$\sigma_q=1$ for $q\equiv3\pmod4$, and $\sigma_q=3$ otherwise. Set
\[
 A_q=(p+\sigma_q q)/4,\qquad R_q=\sigma_q q.
\]
The complete graph has an edge $q\to r$ for every prime $r\mid A_q$ with
$(r/p)=-1$, retaining its weight $v_r(A_q)$. A hard prime has
$(2/p)=(3/p)=(5/p)=(7/p)=1$. Factoring a least positive nonresidue proves the
vertex set is nonempty and its primes are at least 11. Multiplicativity gives
$(A_q/p)=-1$, hence every vertex has an outgoing factor. Since $A_q<p$ and
$q\nmid A_q$, the new vertex is below $p$ and distinct from $q$.

At an original source $a=A_q=\prod_t t^{e_t}$ retain
\[
 \beta\in\prod_t[-e_t,e_t]_{\mathbb Z},\quad
 u=\prod_t t^{e_t+\beta_t},\quad
 v(\beta)=\prod_t t^{\beta_t}\pmod R,\quad
 \mu(g)=\#\{\beta:v(\beta)=g\}.                         \tag{2}
\]
The inverse is $\beta_t=v_t(u)-e_t$. Its total mass is $\tau(a^2)$.
The original two channels are
\[
 E:\ R\mid4u+1\iff v(\beta)=-p^{-1},\qquad
 M:\ R\mid u+a\iff v(\beta)=-1.                         \tag{3}
\]
The third labelled target $-p$ uses the involution $\beta\mapsto-\beta$ before
canonical exterior reconstruction; it is not another channel.

For a hit, $d_0=\gcd(a,u)$ gives
\[
 (h,r,s)=(d_0^2/u,u/d_0,a/d_0),\qquad a=hrs,\ u=hr^2,\ \gcd(r,s)=1.
\]
The ordered returns and inverses are
\[
\begin{array}{c|c|c}
 & (x,y,z)&\text{ordered divisor inverse}\\\hline
 E& (a,hs\kappa,phr\kappa),\quad\kappa=(pr+s)/R
      &u=a^2/(Ry-pa)\\
 M& (a,phs\lambda,phr\lambda),\quad\lambda=(r+s)/R
      &u=pa^2/(Ry-pa).
\end{array}                                                   \tag{4}
\]
Prime valuations give the normalization, and multiplication by $phrs\kappa$ or
$phrs\lambda$ verifies $4/p$. These are inherited Type I/II coordinates
\cite[\S2, Propositions 2.2 and 2.6]{ET}. They also hold at the separately marked
square sources below, where $p/4<a<p$ and $\gcd(pa,R)=1$. If an exterior source
has $a>p/2$, the coordinate switch $(h,r,s,\kappa)\mapsto(h,\kappa,s,r)$ swaps
the first two denominators and returns to the first-half atlas, as proved in
Turn 2 \cite[\S2]{Turn2}. No large auxiliary source is silently deleted.

For a prime residual $q\equiv3\pmod4$, put
\[
 G_q=(\mathbb Z/q\mathbb Z)^\times,\quad
 W_q=\operatorname{supp}\mu_q,\quad
 K_q=\operatorname{Stab}_{G_q}(W_q),\quad
 \Gamma_q=\langle-1,2,t:t\mid A_q\rangle.                 \tag{5}
\]
The full and effective indices are respectively $|G_q/K_q|$ and $|\Gamma_q/K_q|$.
For the triangle below both are six, established rather than assumed.

\begin{lemma}[Inherited primitive form, with its full argument]\label{lem:primitive}
At a failed prime-residual vertex with $|\Gamma_q/K_q|=6$, exactly one simple
prime $r\mid A_q$ lies outside $K_q$, it is the only nonresidue prime factor,
and $A_q=B_qr$ with $W(B_q)=K_q$. Its quotient image has order six, and
$pK_q=rK_q$ or $r^{-1}K_q$.
\end{lemma}
\begin{proof}
Because $-1\notin W_q$ and $K_q\subseteq W_q$, the cyclic group $K_q$ has odd
order; it lies in the quadratic residues. The image of $p$ is a non-square in
$\Gamma_q/K_q\cong C_6$. It cannot be the unique involution, since then an
exterior target is the identity. Hence it has order six, and all three targets
are distinct. The aperiodic quotient support consequently has at most three
points. For every prime outside $K_q$, its signed interval cannot fill its
nontrivial cyclic subgroup, or that period would persist in the final support.
Building each such interval from $2e$ copies of $\{1,g\}$ shows strict growth
at every step. Thus $1+2\sum_{g\ne1}e\le3$. There is at least one active prime,
since $A_q$ is a nonresidue and all elements of $K_q$ are squares. It is unique
and simple; its quotient order is greater than three, hence six. The image of
$p$ is $r$ or $r^{-1}$, rather than the forbidden involution. Each of the three
occupied quotient positions has one active exponent above it, so its fine
support is a translate of $W(B_q)$. Actual $K_q$-periodicity then forces
$W(B_q)=K_q$. The factorwise reciprocity identity identifies the same unique
prime as a nonresidue modulo $p$.
\end{proof}

\section{A closed primitive sextic triangle}
Let
\[
 L_0=526211146410389771,\quad L_1=1165340089472441,\quad
 L_2=23500235136011.
\]
All four integers $p,L_0,L_1,L_2$ are prime. Their exact factorizations give
\[
\begin{array}{c|r|l|c}
 q&A_q& A_q=B_q r_q&r_q\\\hline
31&1877521370392270702928&2^4L_0\cdot223&223\\
223&1877521370392270702976&2^7\cdot41L_1\cdot307&307\\
307&1877521370392270702997&17\cdot19\cdot79\cdot101L_2\cdot31&31.
\end{array}                                                   \tag{6}
\]
\begin{theorem}[Closed minimal-length primitive component]\label{thm:triangle}
At the hard prime (1), the three original vertices form the closed component
\[
 \boxed{31\longrightarrow223\longrightarrow307\longrightarrow31.} \tag{7}
\]
At each vertex $K_q=G_q^6$, $W(B_q)=K_q$, $v_{r_q}(A_q)=1$, and $r_qK_q$
generates $G_q/K_q\cong C_6$. Both complete local channels are empty.
Each vertex has exactly its displayed outgoing nonresidue prime, of weight one.
A canonical cycle consisting of primes $3\pmod4$ cannot have length two, so
three is the minimum possible length. No least-prime assertion is made.
\end{theorem}
\begin{proof}
The primality proof is the finite complete-order certificate explained below;
all factorizations in (6) are exact integer identities. The following bounded
packets already fill the respective sixth-power subgroups:
\[
\begin{array}{c|c|c|c}
q&|K_q|&\text{generator of }K_q&\text{fixed logarithmic packet}\\\hline
31&5&2&[-4,4]\\
223&37&2&[-7,7]+11[-1,1]\\
307&51&171&17[-1,1]+7[-1,1]+[-3,3].
\end{array}                                                   \tag{8}
\]
The first row covers the five residues. The second covers every integer from
$-18$ through $18$. For the third, the actual prime logs of $17,19,79,101$
with respect to $171=4^{10}\pmod{307}$ are $-17,7,-1,-2$. The last two give
$[-3,3]$. The nine centres from the first two are
\[
 -24,-17,-10,-7,0,7,10,17,24.
\]
Consecutive gaps are at most seven, so their radius-three intervals cover
$[-27,27]$, hence all residues modulo 51. This proves all three saturations
by original bounded exponent vectors. Every remaining $L_i$ is in $K_q$:
its residue is respectively $8,136,216$, with logarithm $3,26,39$ relative to
$2,2,4$. Adding those actual factors preserves the already saturated support.

The successor orders in $G_q$ are $6,222,306$; their quotient orders are six.
Also $4\in K_q$ at every vertex, so (6) implies
$pK_q=r_qK_q$. Thus the full support is exactly
\[
 W_q=K_q\{r_q^{-1},1,r_q\}.                              \tag{9}
\]
With $r_qK_q$ labelled 1 in $C_6$, the support classes are $5,0,1$, while the
middle and two exterior target classes are $3,2,4$. The set $\{5,0,1\}$ has
no nontrivial period, proving that $K_q$ is the \emph{actual} stabilizer, not a
guessed power-residue subgroup. The support and successor generate the whole
unit group, so the effective index is also six.

For any prime $t\mid A_q$, reciprocity and $p\equiv-q\pmod t$ give
$(t/p)=(t/q)$; the same identity for $t=2$ follows from the supplementary law.
Every factor of $B_q$ is a square modulo $q$, whereas $r_q$ has odd quotient
exponent and is a nonresidue. Hence the listed successor is the only eligible
prime and appears once. All three vertices themselves are nonresidues modulo
$p$, since $(p/q)=-1$ and $p\equiv1\pmod4$. This proves closure with no omitted
outgoing factor.

For a hypothetical two-cycle $q\leftrightarrow r$ with both primes
$3\pmod4$, both edge laws would give $(r/q)=(q/r)=-1$. Reciprocity makes these
symbols opposite. This proves the minimum-length assertion.
\end{proof}

\paragraph{Primality and original multiplicities.}
The supplied certificate has 80 nodes, of which 42 terminate by complete trial
division below 10,000 and 38 use the following elementary lemma. If the complete
factorization of $n-1$ is into already proved primes and an integer $a$ satisfies
\[
 a^{n-1}\equiv1\pmod n,\qquad
 \gcd(a^{(n-1)/r}-1,n)=1\quad\text{for every }r\mid n-1,
                                                               \tag{10}
\]
then for every prime divisor $l$ of $n$, the order of $a\pmod l$ contains each
full prime-power factor of $n-1$. Consequently $n-1\mid l-1$, so $l\ge n$;
hence $n$ is prime. At the root (1),
\[
 p-1=2^5\cdot3\cdot5\cdot15646011419935589191,
 \qquad a=11.
\]
Every child factor, exponent, base, power and gcd in (10) is in
\texttt{prime\_certificates.json}; both standard-library implementations verify
all 80 nodes. No probable-prime assertion is an input to the final certificate.

The inactive masses are respectively $27,135,243$, so the quotient count vector
in classes $0,1,2,3,4,5$ is
\[
 T_q(1,1,0,0,0,1),\qquad T_q=27,135,243.                 \tag{11}
\]
The complete original divisor counts are $81,405,729$, totalling 1215. Every
vector, integer divisor, fine residue, quotient coordinate and kernel coordinate
is recorded. For example, the inactive counts at $q=31$ on
$1,2,4,8,16$ are $5,6,5,5,6$; they are not $27/5$ per element.

A separate original state demonstrates that (1) is not an ES counterexample:
\[
 R=3,\quad a=1877521370392270702921,
 \quad u=71,\quad (h,r,s)=(71,1,26443962963271418351),
\]
\[
 \kappa=2512176481510784743344,
 \quad (x,y,z)=(a,hs\kappa,ph\kappa).                   \tag{12}
\]
Here $a=19\cdot71\cdot19371487\cdot71847221867$, with all factors separately
certified prime. Equation (4) verifies the positive reciprocal identity.

\section{Higher reciprocity is compatible with the closed cycle}
Let $\omega^2+\omega+1=0$, and retain the primary Eisenstein primes
\[
 \pi_0=-5-6\omega,\qquad \pi_1=-17-6\omega,
 \qquad\pi_2=-17-18\omega.                               \tag{13}
\]
Their norms $a^2-ab+b^2$ are $31,223,307$. Each is $1\pmod3$ in
$\mathbb Z[\omega]$, using the stronger primary convention of Damg\aa rd and
Frandsen \cite[\S2, p.4]{DF}. The reduction maps send $\omega$ to
\[
 z_0=25,\qquad z_1=183,\qquad z_2=289.
\]
For $i\ne j$, define $H_{ij},J_{ij}\in\mathbb Z/3$ by
\[
 \left[\frac{\pi_j}{\pi_i}\right]_3=\omega^{H_{ij}},\qquad
 \left[\frac{\overline\pi_j}{\pi_i}\right]_3=\omega^{J_{ij}}.
\]
The exact tables are
\[
 H=\begin{pmatrix}*&1&0\\1&*&0\\0&0&*\end{pmatrix},\qquad
 J=\begin{pmatrix}*&0&1\\0&*&1\\2&2&*\end{pmatrix},\qquad
 H+J=\begin{pmatrix}*&1&1\\1&*&1\\2&2&*\end{pmatrix}.     \tag{14}
\]
Each entry is checked by reducing the actual Eisenstein numerator at $z_i$,
raising to $(q_i-1)/3$, and comparing with $1,z_i,z_i^2$. Cubic reciprocity gives
$H_{ij}=H_{ji}$; conjugation and reciprocity give $J_{ij}=-J_{ji}$
\cite[\S4, p.8]{DF}. The rational norm character is $H_{ij}+J_{ij}$, not the
primary-only character. The asymmetric conjugate table cannot be discarded.
All forward and backward rational prime characters are nontrivial, consistently
with (7), not in contradiction with reciprocity.

With successor $r_i$ and predecessor $q_{i-1}$, all phases here are positive:
$pK_i=r_iK_i$. Therefore
\[
 q_{i-1}K_i=r_i^{-2}K_i,\qquad
 q_{i-1}r_i^2\pmod{q_i}=16,196,17.                       \tag{15}
\]
These kernel elements have explicit sixth roots $3,38,20$, respectively.
This is the exact neighbour-square recurrence of the predecessor, now satisfied
around a closed primitive triangle \cite[\S3]{BryanChain}.

The section $j\mapsto r_i^j$ has carry
\[
 \omega_i(j,k)=(r_i^6)^{\lfloor(j+k)/6\rfloor},\qquad
 r_i^6\pmod{q_i}=1,8,144.                                \tag{16}
\]
Every fine group element has coordinates $(j,k_0)$ with $k_0\in K_i$; the product
uses that carry, and its inverse is $g\mapsto(j,gr_i^{-j})$. At $q=307$ all lifts
of a quotient generator have order divisible by 18. The extension of $C_6$ by
$C_{51}$ does not split; $C_{306}$ must not be replaced by $C_6\times C_{51}$.
The quotient labels at different vertices are separately marked by their own
successor. Their use of a common abstract $C_6$ is not a residue-group
homomorphism between different prime moduli.

\section{The smallest odd coefficient component, including its nilpotent part}
The following comparison is an explicit classical algebraic CRT calculation.
It states exactly what the primitive packet does in the odd component; it is
not proposed as an independent ES existence theorem.
Put $\varepsilon^2=0$ and $\mathbb F_4=\mathbb F_2[\omega]/(\omega^2+\omega+1)$.
There is a unital algebra isomorphism
\[
 \boxed{\Phi:\mathbb F_2[X]/(X^6-1)\longrightarrow
 \mathbb F_2[\varepsilon]/\varepsilon^2\times
 \mathbb F_4[\varepsilon]/\varepsilon^2,\quad
 F\mapsto(F(1+\varepsilon),F(\omega(1+\varepsilon))).}    \tag{17}
\]
Its complete inverse can be written without solving a linear system. Set
\[
 e_0=1+X^2+X^4,\quad e_1=X^2+X^4,\quad E=1+X^3,\quad Z=X^4.
\]
For $(a+b\varepsilon,c+d\varepsilon)$, with $a,b\in\mathbb F_2$ and
$c,d\in\mathbb F_4$, the inverse is
\[
 \boxed{(a+bE)e_0+\bigl(c(Z)+E d(Z)\bigr)e_1.}           \tag{18}
\]
Indeed $e_0,e_1$ are complementary orthogonal idempotents; $E^2=0$; $Z$ on the
second component is $\omega$; and $e_0X=e_0(1+E)$,
$e_1X=e_1Z(1+E)$. These identities prove both compositions and multiplicativity.
Equivalently, $X^6-1=(X+1)^2(X^2+X+1)^2$ with coprime factors.

Every original integer $T_q$ in (11) is odd. Its reduction is the same marked
polynomial $P=1+X+X^5$, and direct substitution gives
\[
 \boxed{\Phi(P)=(1,\varepsilon),\qquad
 \Phi(P^2)=(1,0),\qquad P^2=e_0.}                        \tag{19}
\]
The cubic body is zero while its cubic nilpotent coefficient is nonzero. In the
source convention $G(A)=\{\tau\}\sqcup A^\bullet$, this is a supported zero
observation with a retained jet, not absence. It does not force a target:
the fine, complete sources in Theorem~\ref{thm:triangle} have all three target
classes absent despite their nonzero cubic jet.

Squaring $P$ introduces two independent active occurrences and their pair
multiplicities. The original active prime has exponent bound one. Neither that
squared source nor its exterior target coefficients are an available move inside
the original box. All integer counts and fine kernel fibres preceding reduction
remain in the certificate. Equality of the three reduced local packets is not
an invented isomorphism of their original factor sources.

\section{A genuinely additional source breaks the minimal cycle}
\begin{theorem}[Square-nine escape, with no new factor-supply assumption]
\label{thm:nine}
Let $q_0\to q_1\to q_2\to q_0$ be any canonical directed triangle of external
nonresidue primes, all $3\pmod4$, at a hard prime $p$. It need not be primitive
and its vertices may have other canonical exits. For every vertex $q_i$, the
new original source
\[
 A_{q_i,9}=(p+9q_i)/4
\]
is below $p$, has quadratic character $-1\pmod p$, and has at least one
nonresidue prime factor. \emph{Every} such factor lies outside the triangle.
\end{theorem}
\begin{proof}
Let $Q$ be the largest vertex. Its incoming edge gives
$4cQ=p+q_{\rm prev}<p+Q$, so $3Q<p$. Thus $9q_i<3p$ and the new sources are
positive integers below $p$. Multiplication by the square 9 preserves the
nonresidue character, proving existence of a nonresidue factor below $p$.
No source is divisible by its own vertex. Further,
\[
 \gcd(A_{q,1},A_{q,9})=\gcd(A_{q,1},2),                 \tag{20}
\]
so its canonical successor, an odd prime at least 11, cannot divide $A_{q,9}$.
Finally its predecessor $s$ satisfies $(q/s)=-1$. Both are $3\pmod4$, hence
$(s/q)=1$. Any nonresidue factor $t$ of $A_{q,9}$ has $(t/q)=-1$: reciprocity
in $p\equiv-9q\pmod t$ gives $(t/p)=(t/q)$. The predecessor is therefore
excluded too. These are all three vertices.
\end{proof}

The arithmetic map in this theorem changes the source, retaining
$(p,q,1,A_{q,1})$ and $(p,q,9,A_{q,9})$ separately. It is not a transport of an
old missing divisor or an enlargement of its multiplicity. Both have their
original factorization, exponent box, channels and return (4).

\begin{theorem}[Square-source cardinality obstruction]\label{thm:ports}
Let $C$ be a finite set of external nonresidue primes, all $3\pmod4$, and let
$L\ge1$. Suppose
\[
 (2L-1)^2\max C<3p,\qquad \min C>2L-2.                  \tag{21}
\]
If every nonresidue factor of every source
$A_{q,j}=(p+(2j+1)^2q)/4$, $q\in C$, $0\le j<L$, lies in $C$, then
\[
 \boxed{|C|\ge2L+1.}                                   \tag{22}
\]
\end{theorem}
\begin{proof}
Every source is an actual integer below $p$, with nonresidue character, hence
has a nonempty outgoing factor set. For $i<j$,
\[
 \boxed{\gcd(A_{q,i},A_{q,j})
 =\gcd(A_{q,i},(j-i)(i+j+1)).}                          \tag{23}
\]
This follows by subtracting the sources and cancelling their unit $q$. Every
prime divisor of the two factors on the right is at most $2L-2$. By (21), an
internal prime cannot occur at two different ports. Each vertex therefore has
at least $L$ distinct internal successors. There are no self-edges, and no
opposite pair of edges: the square-source character law would otherwise make
both reciprocal symbols $-1$ at two $3\pmod4$ primes. Hence the total number
of directed edges is at most $|C|(|C|-1)/2$ and at least $L|C|$. This proves (22).
\end{proof}

\begin{corollary}[Uniform escape for small primitive components]
Let $C$ be a closed canonical strongly connected component consisting entirely
of primitive effective-index-six prime-residual failures. If $|C|\le6$, at least
one of its square sources with multiplier 9 or 25 has a nonresidue prime factor
outside $C$. No assertion that this factor's next selector is occupied is made.
\end{corollary}
\begin{proof}
Turn 1's actual normal form gives $A_q=B_qr_q$, $W(B_q)=K_q$, and a unique
outgoing prime $r_q$ \cite[\S5]{Turn1}. Each $B_q\ge4$. If $B_q=1$, then
$K_q=1$ and $|\Gamma_q|=6$; since $2\in\Gamma_q$, $q\mid2^6-1=63$, impossible
for a hard external prime. If $B_q=2$, its three-element signed support can be
a subgroup only if $\operatorname{ord}_q(2)=3$, forcing $q=7$. If $B_q=3$, the
same argument forces $q=13$, contrary to $q\equiv3\pmod4$.
At a largest vertex $Q$, the incoming canonical edge now gives
$16Q\le p+q_{\rm prev}<p+Q$. Thus $15Q<p$ and $25Q<3p$.
Take $L=3$ in Theorem~\ref{thm:ports}; $\min C\ge11>4$. Closure under all three
ports would require at least seven vertices. Port 1 is already closed, so an
outside factor occurs at 9 or 25.
\end{proof}
For an arbitrary all-$3\pmod4$ canonical closed component one still has
$3\max C<p$. The same argument with $L=2$ proves that components of at most
four vertices cannot remain closed after multiplier 9 is added.

\section{Actual escape, exact scope, and the next obligation}
For (1), the three multiplier-nine factorizations are
\[
\begin{array}{c|l|l}
q&A_{q,9}&\text{all nonresidue prime factors}\\\hline
31&2\cdot5\cdot11\cdot131\cdot2467\cdot52814328044617&11\\
223&2\cdot7\cdot17\cdot69623\cdot113306597729903&113306597729903\\
307&7\cdot268217338627467243373&268217338627467243373.
\end{array}                                                   \tag{24}
\]
All factors in the table are certified prime. At $q=31$ the full new source has
four E and four M states. One M state is
\[
\begin{gathered}
 R=279,\quad a=1877521370392270702990,\quad u=156417668,\\
 (h,r,s,\lambda)=(323177,22,264071640223085,946493334133),
\end{gathered}                                                \tag{25}
\]
with ordered denominators given by (4). At $q=223$ and $307$, both new channels
are empty. Thus even the explicit escape does not mean that every added source
is immediately occupied. At $q=223$, projection from $9q$ to $q$ supplies two
candidates in each channel, but none lifts with its original exponent vector.
For instance,
\[
 u=1098476221603538346086974,\qquad
 223\mid4u+1,\quad u\equiv4\pmod9,
\]
whereas E requires $u\equiv2\pmod9$. The other projected E candidate is
$5\pmod9$, and both projected M candidates are $1\pmod9$ instead of the
required $8\pmod9$. The square factor of the residual remains part of the
actual selector; it is not discarded because it has trivial Jacobi sign.

The actual closed triangle refutes a strict decrease on every canonical primitive
edge, even with all cubic and nilpotent data retained. Square sources provide a
proved extra arithmetic input, but escape is not eventual selector occupancy.
Larger or mixed-type closed components remain possible. Turn 4 must retain the
size/type bounds (21)--(22), or pivot to the complete-shell signed coefficient sum.

\paragraph{Exact computational scope.}
The standard-library checks cover all 1215 original vectors, all fine products
and carries, both square-port gates, all 80 prime-proof nodes, and the complete
64-element coefficient ring. The separate triangle scan covers the 636 hard
primes up to 250,000: 1,802,376 vertices, 961,554 edges, and 19 triangles, all
obeying square-nine escape. It is not an all-shell ES census or the large prime's
full graph. Independent code, hashes and negative controls accompany the proof.
No least-prime, historical-priority, independent-review, Lean-build or universal
ES claim is made.
